\section{Red-team: the strongest objections, answered}
\label{sec:redteam}

A specification that cannot survive its own strongest critics is marketing.
We state the seven attacks a hostile architect would mount and answer each on
the record. Where an objection lands, we concede the exact scope.

\begin{redteambox}[O1. ``This is just a neuromorphic chip with extra steps.'']
Neuromorphic machines co-locate memory and compute and are event-driven ---
which \TALOS{} also is --- but their topology is \emph{ad hoc}, they have no
native error correction, and their state object is a heuristic (a spike
train), not a derived one. The distinguishing content of \TALOS{} is
Theorem~\ref{thm:collapse}: the topology, the code, the instruction algebra,
and the learning rule are \emph{the same object}, so the fabric is provably a
distance-3 code (\S\ref{sec:ecc}) and the routing is provably third-order
(\S\ref{sec:fabric}). No neuromorphic design has, or can bolt on, a topology
that is simultaneously its ECC and its instruction algebra without being the
Fano plane --- which forces $N=7$ (Theorem~\ref{thm:unique}). \emph{Verdict:
overlaps in flavour, differs in the load-bearing theorem.}
\end{redteambox}

\begin{redteambox}[O2. ``$\Dens(\mathbb{C}^7)$ is quantum --- you need cooling and it will decohere.'']
No. UHM's substrate-closure theorem (T-153) allows a \emph{classical} digital
substrate to realise the same $\Gam$: the complex entries are algebraic
floating-point state variables ($48$ real dynamical parameters, $784$~bytes of
\texttt{f64} in L1 cache at room temperature, carrying the dual-aspect phase
Gap, T-132), not a physical macroscopic quantum superposition of cryogenic
pointer states. ``Density matrix'' names the mathematical type (Hermitian, PSD,
unit trace) enforced on a classical register, not a cooling requirement.
\emph{Verdict: category error; conceded nowhere.}
\end{redteambox}

\begin{redteambox}[O3. ``BQP-bounded means no computational advantage --- why bother?'']
We agree on the premise and reject the inference. UHM proves $\Regen$ grants
no speed-up beyond BQP, and we foreground this in \S\ref{sec:intro}. The
advantage claimed is never complexity class. It is (i)~a tiny, L1-resident
$784$-byte state removing off-chip DRAM fetches per tick, together with a
logical thermodynamic decomposition isolating entropy export to adaptation and
maintenance (\S\ref{sec:energy}); (ii)~sector-level distance-$3$ Hamming
duality built directly into the seven-line fabric without a separate routing
table (\S\ref{sec:ecc}); (iii)~native, single-loop anchored adaptation
(\S\ref{sec:selfmodel}); and (iv)~exact Fano pinching contraction
$\Phi\to\Phi/9$ and $\mathrm{PSL}(2,7)\ltimes U(1)^7$ fabric uniformity
(\S\ref{sec:fabric}). None of these is a complexity-class claim, and all are
testable (\S\ref{sec:benchmarks}). \emph{Verdict: the objection refutes a
claim we do not make.}
\end{redteambox}

\begin{redteambox}[O4. ``Seven is numerology dressed as engineering.'']
Seven is forced without a tie-break by binary projective coding alone.
Theorem~\ref{thm:unique} proves that requiring the binary incidence matrix $M$
of $\mathrm{PG}(2,q)$ to generate a perfect binary Hamming $[2^r{-}1,2^r{-}1{-}r,3]$
code already excludes \emph{every} odd prime power $q$ --- including the bare
integer coincidence $q=5, N=31$, where
$\operatorname{rank}_{\mathbb{F}_2}(M_{\mathrm{PG}(2,5)})=30$ leaves only a
one-dimensional kernel (verified in Lean~4,
\texttt{pg2\_5\_refutes\_N31\_hamming\_fusion}, and in
\texttt{verify\_foundational\_bounds.py}) --- and every even prime power
$q=2^m$ ($m\ge 2$) modulo~$4$ ($q=90, N=8191$ is not even a prime power).
Binary plane--code fusion selects $N=7$ uniquely across all prime powers
$q\ge 2$, and Hurwitz's unique non-associative division algebra $\Oct$
independently lands on the exact same integer $7$. \emph{Verdict: answered
directly by machine-checked Theorem~\ref{thm:unique}.}
\end{redteambox}

\begin{redteambox}[O5. ``Substrate closure says any substrate works --- so the architecture is not special.'']
Correct and irrelevant to the claim. T-153 says any substrate meeting
C1--C3 \emph{runs} the logical architecture; it does not say all substrates
run it \emph{equally well}. \TALOS{} is the impedance-matched member of the
family (\S\ref{sec:realizability}): the one whose physical primitives are the
logical primitives, so it approaches the energy floor a general-purpose
substrate cannot. The uniqueness claim is about the \emph{logical}
architecture (Theorem~\ref{thm:unique}); the multiplicity is about
\emph{physical} substrates. Both are theorems, and they are compatible.
\emph{Verdict: conceded and reframed; no contradiction.}
\end{redteambox}

\begin{redteambox}[O6. ``The energy floor is theoretical --- real chips never reach Landauer.'']
Granted, and stated explicitly in \S\ref{sec:energy}: no CMOS chip approaches
$k_BT\ln 2$, and reversible Landauer bounds apply equally to \emph{any}
logically reversible architecture (including Bennett-reversible von~Neumann
circuits). We therefore do not compare Landauer floors against physical DRAM
joules. The defensible engineering claims are narrower and rung-dependent: on
digital chips, keeping a $784$-byte state resident in L1/registers removes
off-chip DRAM fetches per tick; on future physically reversible or
driven-dissipative substrates, the unitary/non-unitary tick split isolates
logical entropy export ($\le\log_2 7\approx 2.807$~bits per state) to
maintenance and learning.
\emph{Verdict: conceded scope; no apples-to-oranges Landauer claim is made.}
\end{redteambox}

\begin{redteambox}[O7. ``Unbuildable --- show me one that runs.'']
The bottom rung runs today: \SYNARC{} holds \DM{} in L1, advances it with a
Fano-sparse Strang kernel, and reproduces every UHM invariant bit-for-bit
against an independent oracle (M1 passed). That realisation emulates the
fabric in software and therefore forgoes the energy floor --- honestly, the
today-machine buys correctness, not efficiency. The efficiency claims are
rung-dependent (\S\ref{sec:realizability}) and are offered as a roadmap with
falsifiable milestones (\S\ref{sec:benchmarks}), not as a shipping product.
\emph{Verdict: the logical architecture is demonstrated; the efficient
substrate is a staged research program, labelled as such.}
\end{redteambox}

\noindent The two objections that land --- O6 and O7 --- both concern the gap
between the logical architecture (proven, and running) and the
impedance-matched substrate (a research ladder). We concede that gap
explicitly and it is the subject of the roadmap (\S\ref{sec:discussion}).
None of the seven touches the two theorems of \S\ref{sec:collapse}.
